% ECONOMICS_PROBLEM: {"id": "C78-417", "title": "Stability probability under compact indifference", "jel_code": "C78", "source_id": "OP-0544"}
% FIRST_POSTED: 2026-10-10
% ATTEMPT_STATUS: solved
\documentclass[11pt]{article}
\usepackage[margin=0.85in]{geometry}
\usepackage{amsmath,amssymb,amsthm}
\usepackage[T1]{fontenc}
\usepackage[hidelinks]{hyperref}
\hypersetup{pdftitle={Counting stable tie refinements},pdfauthor={Ji Ho Bae}}
\newtheorem{theorem}{Theorem}
\newtheorem{lemma}{Lemma}
\newcommand{\FP}{\mathsf{FP}}
\newcommand{\PP}{\mathsf{PP}}
\newcommand{\sharpP}{\mathsf{\#P}}
\newcommand{\Stab}{\operatorname{Stab}}
\newcommand{\PM}{\operatorname{PM}}
\setlength{\parindent}{0pt}
\setlength{\parskip}{4pt}
% AI assistance: OpenAI Codex with automated collaboration on 10 October 2026.
% Model version was not exposed. No human edits to the mathematical proof.
% Citation, provenance, and reproduction metadata are recorded in CitationAudit.json.
\begin{document}
\begin{center}
{\Large\bfseries Counting stable tie refinements}\\[5pt]
{\small Ji Ho Bae}
\end{center}

There are equally many agents on two sides. Each agent gives a complete
weak preference order, and $M$ is a specified perfect matching.
Break each tie uniformly and independently into a strict order.
An unmatched pair blocks $M$ when both agents prefer each other to their
partners.
This is the compact-indifference model of
\cite[Section 3.1.2, Table 1]{survey} and
\cite[Sections 2 and 5, Table 3]{uncertain}.
If the tie sizes are $t_{a,j}$, set
\[
 D(I)=\prod_{a,j}t_{a,j}!,\qquad
 N(I,M)=\#\{\text{strict refinement profiles in which }M\text{ is stable}\}.
\]
Thus $p_I(M)=N(I,M)/D(I)$. The weak orders are explicit and any rational
threshold is encoded in binary.

\begin{theorem}\label{thm:classification}
The integer function $N$ is $\sharpP$-complete under polynomial-time
Turing reductions. Computing the exact rational probability $p$ and
computing $N$ are polynomial-time Turing equivalent.
The language
\[
 \{(I,M,\tau):\ \tau\in[0,1]\text{ is rational and }p_I(M)\geq\tau\}
\]
is $\PP$-complete under polynomial-time Turing reductions.
No complexity separation is assumed. Here Turing reductions allow
polynomially many adaptive oracle queries.
\end{theorem}

Counting perfect matchings of a simple balanced bipartite graph is equivalent
to computing the permanent of a $0$--$1$ matrix, which is
$\sharpP$-complete~\cite[Theorem 1]{valiant}. Section 2 of that paper
uses polynomial-time oracle reductions.

For a simple bipartite graph $H=(V,E)$ of minimum degree at least two,
write $h=|V|$, $e=|E|$, and $d_v=\deg_H(v)$.
Let $C(H)$ be the number of ways for every vertex to choose one incident
edge such that no edge is chosen by both endpoints.

\begin{lemma}\label{lem:inner}
$C(H)$ can be computed in polynomial time with an exact
compact-indifference stability-probability oracle.
\end{lemma}
\begin{proof}
If $H$ is empty, $C(H)=1$. For nonempty $H$ and an integer
$t\geq\Delta(H)$, give each original vertex $v$ a unique
private partner on the opposite side. Add $r_v=t-d_v$ distinct probe
agents on the opposite side, each with its own private partner.
The matching $M$ consists of these private pairs.
Every original vertex has one top tie containing its partner, its
$H$-neighbors, and its probes; it has size $t+1$.
A probe of $v$ strictly ranks $v$ first, its own partner second, and all
others below. Every private partner ranks its matched agent first.
Put all remaining names in a fixed strict order below the partner.
Every private pair adds one agent to each side, so the sides are balanced.

Only edges of $H$ and original-vertex--probe pairs can block: every other
unmatched pair has an endpoint that ranks its partner above the other
agent. Generate each top-tie permutation using independent uniform
priorities in $[0,1]$, smaller first.
Write $x_v$ for its partner's priority.
Conditional on all $x_v$, a given tied competitor precedes the partner
with probability $x_v$. Different unmatched pairs use distinct competitor
priorities, so their safety events are conditionally independent.
A probe pair is safe with probability $1-x_v$; an edge $vw$ of $H$
is safe with probability $1-x_vx_w$. Therefore the oracle gives
\[
 p_H(t)=\int_{[0,1]^h}
       \prod_v(1-x_v)^{t-d_v}
       \prod_{vw\in E}(1-x_vx_w)\,dx.               \tag{1}
\]
The $h$ top ties are the only nontrivial ties, giving
$D=((t+1)!)^h$ refinement profiles.

For $F\subseteq E$, let $k_v=\deg_F(v)$.
Expanding the finite edge product and using
\[
 \int_0^1 x^k(1-x)^r\,dx=\frac{k!\,r!}{(r+k+1)!}
\]
gives, for the nonnegative integers $r=t-d_v$,
\[
 Q(t):=p_H(t)\prod_v(t-d_v+1)
      =\sum_{F\subseteq E}(-1)^{|F|}
          \prod_v L_{t-d_v}(k_v),\qquad
 L_r(k)=\frac{k!}{\prod_{j=2}^{k+1}(r+j)}.         \tag{2}
\]
The empty denominator product is one. The integral identity follows by
$k$ integrations by parts, starting with
$\int_0^1(1-x)^r dx=1/(r+1)$.
Set
\[
 P(t)=\prod_v P_v(t),\qquad
 P_v(t)=\prod_{j=2}^{d_v+1}(t-d_v+j),\qquad
 K=\sum_v d_v=2e.
\]
Multiplying each term in (2) by $P$ cancels all denominators:
\[
 A(t):=P(t)Q(t)
   =\sum_{F\subseteq E}(-1)^{|F|}
       \prod_v\left[
          k_v!\prod_{j=k_v+2}^{d_v+1}(t-d_v+j)
       \right].                                  \tag{3}
\]
Thus $A$ is a polynomial of degree at most $K$.
Query (1) at the $K+1$ valid integers
$t=\Delta(H),\ldots,\Delta(H)+K$, calculate $A(t)$, and interpolate
its value at zero. All queries use nonnegative probe counts; evaluation
at zero takes place only in polynomial (3).

Since $d_v\geq2$, the vertex factor in (3) at zero is
\[
 k!\prod_{j=k+2}^{d_v+1}(j-d_v)
   =
   \begin{cases}
   0,&k\leq d_v-2,\\
   (d_v-1)!,&k=d_v-1,\\
   d_v!,&k=d_v.
   \end{cases}
   =(d_v-1)!\binom{k}{d_v-1}.
\]
Consequently
\[
 \frac{A(0)}{\prod_v(d_v-1)!}
      =\sum_{F\subseteq E}(-1)^{|F|}
          \prod_v\binom{\deg_F(v)}{d_v-1}.          \tag{4}
\]
For each $v$, the binomial factor counts a selection $B_v$ of $d_v-1$
incident edges contained in $F$. Exchange the finite sums: a fixed
family $(B_v)$ contributes
$\sum_{F\supseteq B}(-1)^{|F|}$, where $B=\bigcup_v B_v$.
This sum is zero if $B\ne E$, by pairing subsets according to any edge
of $E\setminus B$; it equals $(-1)^e$ if $B=E$.
Selecting $B_v$ omits exactly one incident edge at $v$.
The union is $E$ precisely when no edge is omitted at both endpoints.
The omitted edges therefore give precisely the choices counted by $C(H)$:
\[
 C(H)=\frac{(-1)^e A(0)}{\prod_v(d_v-1)!}.          \tag{5}
\]

Put $T=\Delta+K$. Each query uses $O(h(T+1))$ agents with explicit
complete lists, and its profile denominator has $O(hT\log(T+1))$ bits.
At a query point, $0\leq p_H(t)\leq1$ and all $K+h$ factors in
$P(t)\prod_v(t-d_v+1)$ are at most $T+1$; hence the integer $A(t)$
has at most $O((K+h)\log(T+1))$ bits.
Each Lagrange coefficient at zero uses $K$ numerators of magnitude at
most $T$ and $K$ nonzero denominator differences of magnitude at most
$K$. A common denominator over the $K+1$ summands therefore has
$O(K^2\log(T+1))$ bits. All arithmetic is consequently polynomial.
\end{proof}

\begin{lemma}\label{lem:outer}
The number of perfect matchings of any simple balanced bipartite graph can be
computed in polynomial time with an oracle for $C(H)$ on simple bipartite
graphs of minimum degree at least two.
\end{lemma}
\begin{proof}
Let $G$ be simple and balanced bipartite, with $V=2n>0$ vertices and
degrees $d_v$. The empty graph on
zero vertices has one perfect matching and needs no oracle.
For each integer
$R=\Delta(G)+1,\ldots,\Delta(G)+V+1$, attach
$\ell_v=3(R-d_v)$ disjoint four-cycles to $v$, using one bridge from
$v$ to a designated vertex $z$ of each cycle. Call the resulting graph
$H_R$ and write $L=\sum_v\ell_v$.
Each cycle can be bipartitioned with $z$ opposite $v$, so $H_R$ is
bipartite and simple. The original vertices have degree at least three,
the designated cycle vertices have degree three, and the others degree
two.

Consider one cycle $z-a-b-c-z$ and its bridge to $v$.
If $v$ chooses the bridge, $z$ cannot choose it.
The four cycle vertices must then choose cycle edges with no mutual
choice. Following choices reaches a directed cycle; with mutual choices
excluded it must use all four vertices. There are two orientations.
If $v$ does not choose the bridge, those same two choices remain, and
there are four choices with $z$ choosing $v$.
For the latter count, $b$ chooses either $a$ or $c$; its chosen neighbor
must choose $z$, while the other neighbor has two choices. These two
cases give four.

Factor out six choices per cycle. At an original vertex, choosing an
original edge has weight one, and choosing one of its bridges has
weight $2/6=1/3$. Its total choice weight is therefore
$d_v+\ell_v/3=R$.
Apply inclusion-exclusion to the remaining forbidden events, in which
both endpoints choose an original edge. An intersection
of events sharing an endpoint is impossible, since a vertex chooses
only one edge. The nonzero intersections therefore correspond to
matchings $F$ of $G$. Each fixes choices at $2|F|$ vertices,
with weight one, and leaves total weight $R$ at every other vertex.
We obtain
\[
 Z_G(R):=\frac{C(H_R)}{6^L}
       =\sum_{\substack{F\subseteq E(G)\\F\ \mathrm{matching}}}
          (-1)^{|F|}R^{V-2|F|}.                   \tag{6}
\]
The right side is a polynomial of degree $V$.
Its constant coefficient is $(-1)^n\PM(G)$, since the only matchings
with $V-2|F|=0$ are perfect matchings. Interpolating at the $V+1$
chosen $R$ values therefore recovers $\PM(G)$.

Here $R=O(V)$, $L=O(V^2)$, and $H_R$ has $O(V^2)$ vertices and edges.
The factor $6^L$ has polynomial bit length, and $0\leq Z_G(R)\leq R^V$
by the weighted-choice interpretation. Exact interpolation has the same
polynomial bit bound as above. Combining with Lemma~\ref{lem:inner},
the largest matching instance has $O(V^4)$ agents and $O(V^8)$
preference-list entries, and the total number of queries is polynomial.
\end{proof}

\begin{proof}[Proof of Theorem~\ref{thm:classification}]
Encode every tie permutation by a fixed-length list of its element
indices, rejecting invalid lists. Every complete strict refinement has
exactly one valid encoding. A nondeterministic machine guesses the
concatenation, checks validity, and checks all unmatched pairs for
blocking in polynomial time. Its accepting-branch count is exactly $N$,
so $N\in\sharpP$.

The product $D$ is computable in polynomial time and has polynomial
bit length. Given $N$, compute $p=N/D$; given exact $p$, compute the
integer $N=Dp$. Thus the functions are Turing equivalent. Lemmas
\ref{lem:inner}--\ref{lem:outer} and Valiant's theorem prove
$\sharpP$-hardness and hence the stated completeness of $N$.

For threshold membership, write $\tau=a/b$ with $b>0$ and set
$k=\lceil aD/b\rceil$. Then $p\geq\tau$ is equivalent to $N\geq k$.
Let $B$ be the length of the fixed encoding just used, so the counting
machine has $2^B$ branches and $N\leq D\leq2^B$.
If $k=0$, accept deterministically.
Otherwise construct a machine with $2^{B+1}$ branches by guessing a
leading bit and then $B$ further bits. In its first half it performs the
counting test, giving $N$ accepting branches.
In its second half it interprets the $B$ bits as an integer and accepts
exactly the first $2^B-k+1$ integers. Its accepting count is
$N+2^B-k+1$, strictly above half of all branches exactly when $N\geq k$.
The bound $1\leq k\leq D\leq2^B$ makes that second-half count valid.
The machine runs in polynomial time, proving membership in $\PP$.

A threshold oracle recovers $N$ by binary search over $0,\ldots,D$,
using thresholds $k/D$ and $O(\log(D+1))$ calls.
Thus the threshold oracle computes every $\sharpP$ function.
For any language in $\PP$, take a polynomial-time majority machine
making a fixed $q(|x|)\geq1$ binary guesses.
Its accepting count is a $\sharpP$ function; test whether it exceeds
$2^{q(|x|)-1}$. This gives a polynomial-time Turing
reduction from every $\PP$ language to the threshold problem and proves
the claimed $\PP$-hardness.
\end{proof}

\paragraph{Source and verification.}
\href{https://github.com/jbaelaw/compact-indifference-stability/blob/10b2dac3c6addbd32205f328c222b4c20a48b52b/paper/Proof.pdf}{Four-page proof PDF}
and \href{https://github.com/jbaelaw/compact-indifference-stability/blob/10b2dac3c6addbd32205f328c222b4c20a48b52b/paper/Proof.tex}{complete original generated source}.
The \href{https://github.com/jbaelaw/compact-indifference-stability/tree/10b2dac3c6addbd32205f328c222b4c20a48b52b/lean}{pinned Lean companion}
checks the finite cancellation and polynomial-factor identities;
it does not formalize the whole probability construction or the complexity classes.
Lean 4.35.0-rc3 with mathlib commit
\href{https://github.com/leanprover-community/mathlib4/tree/b84a70d6a5ed793cc46184160f4d4188d8c825fb}{b84a70d6}
passed \texttt{lake build} and the axiom audit without warnings or placeholders.
\href{https://github.com/jbaelaw/compact-indifference-stability/blob/10b2dac3c6addbd32205f328c222b4c20a48b52b/verification/LeanVerification.json}{Verification scope and source hashes}
and \href{https://github.com/jbaelaw/compact-indifference-stability/blob/10b2dac3c6addbd32205f328c222b4c20a48b52b/paper/CitationAudit.json}{checked citation locations}
are recorded separately. The exact-arithmetic tests and complete-profile enumeration
are reproducible from the project. Model: GPT Codex, serving revision not exposed.
The proof was generated with automated collaboration following the user's instructions.
This is a proof claim submitted for expert review.

\section*{Completion Estimate}
100\%
The claimed classification uses polynomial-time Turing reductions, as stated in the theorem.

\begin{thebibliography}{3}
\small
\bibitem{survey}
K. Cechl\'arov\'a, \'A. Cseh, and D. F. Manlove.
\emph{Selected open problems in Matching Under Preferences}.
Bulletin of the EATCS, no.~128 (2019).
Section 3.1.2 and Table 1, PDF pp.~8--9.
\url{https://hpi.de/friedrich/docs/publications/2019/BEATCS.pdf}.

\bibitem{uncertain}
H. Aziz, P. Bir\'o, S. Gaspers, R. de Haan, N. Mattei, and B. Rastegari.
\emph{Stable Matching with Uncertain Linear Preferences}.
Algorithmica 82, 1410--1433 (2020).
Section 2; Section 5; Table 3.
\href{https://doi.org/10.1007/s00453-019-00650-0}{doi:10.1007/s00453-019-00650-0}.

\bibitem{valiant}
L. G. Valiant. \emph{The Complexity of Computing the Permanent}.
Theoretical Computer Science 8, 189--201 (1979).
Section 2, p.~191 (oracle reductions); Theorem 1, p.~194.
\href{https://doi.org/10.1016/0304-3975(79)90044-6}{doi:10.1016/0304-3975(79)90044-6}.
\href{https://www.math.cmu.edu/users/af1p/Teaching/MCC17/Papers/permanent.pdf}{Primary paper scan}.
\end{thebibliography}
\end{document}
